\documentclass[11pt,a4paper]{article}
\usepackage[T1]{fontenc}
\usepackage[utf8]{inputenc}
\usepackage{lmodern}
\usepackage[margin=2.7cm]{geometry}
\usepackage{microtype}
\usepackage{amsmath,amssymb,amsthm}
\usepackage{booktabs}
\usepackage{xcolor}
\usepackage{xurl}
\usepackage[
  colorlinks=true,
  linkcolor={blue!55!black},
  citecolor={green!40!black},
  urlcolor={blue!65!black},
  pdftitle={Small quasigroups and the implication semilattice of Schroeder's 990 laws},
  pdfauthor={Carlo Perassi},
  pdfkeywords={quasigroups, Latin squares, equational theories, varieties, finite models},
]{hyperref}
\usepackage[capitalise,noabbrev]{cleveref}
\urlstyle{same}

\newcommand{\ldiv}{\mathbin{/\!/}}
\newcommand{\rdiv}{\mathbin{\backslash\!\backslash}}
\newcommand{\Sch}[1]{\textsf{Sch-#1}}
\newcommand{\arxiv}[1]{\href{https://arxiv.org/abs/#1}{arXiv:#1}}
\newcommand{\oeis}[1]{\href{https://oeis.org/#1}{#1}}
\newcommand{\doi}[1]{\href{https://doi.org/#1}{doi:#1}}
\newcommand{\repo}{\url{https://github.com/carlok/quadratula}}
\input{numbers.tex}

\newtheorem{lemma}{Lemma}
\theoremstyle{definition}
\newtheorem{definition}{Definition}

\title{Small quasigroups and the implication semilattice\\ of Schröder's 990 laws}
\author{Carlo Perassi}
\date{Draft, \today}

\begin{document}
\maketitle

\begin{abstract}
Le Floch~\cite{lefloch} determined all implications among Schröder's \NLaws{}
quasigroup laws: \NClasses{} equivalence classes and \NVarieties{} closed sets. We
measure how much of this structure is witnessed by quasigroups of small order.
Quasigroups of order at most 4 witness \ClassSatFour{} of the non-implications
between classes, and those of order at most 6 witness \ClassSatSix{} and realise
\VarSix{} of the \NVarieties{} closed sets. Exhaustive search over the models of the
remaining laws and closed sets extends the picture exactly: quasigroups of order at
most \CompleteUpTo{} witness \ClassSatNine{} of the non-implications and realise
\VarNine{} closed sets, with least witness orders up to \WitTwentyThree. A set of
\FullCover{} quasigroups witnesses every non-implication between classes and, given
Le Floch's list of closed sets, no smaller set of quasigroups of any order does.
Code and data: \href{https://github.com/carlok/quadratula}{github.com/carlok/quadratula}.
\end{abstract}

\noindent\textbf{Keywords:} quasigroups, Latin squares, equational laws, varieties,
finite models.\\
\textbf{MSC 2020:} 20N05, 08B05.

\section{Introduction}

Schröder~\cite{schroeder} considered the equations with three variables and two
operations on each side, such as $x \circ_1 (y \circ_2 z) = \alpha \circ_3 (\beta
\circ_4 \gamma)$, where each $\circ_i$ is independently one of the quasigroup
operations $*, \ldiv, \rdiv$, either side may instead be bracketed as
$(\cdot \circ \cdot) \circ \cdot$, and $(\alpha,\beta,\gamma)$ is a permutation of
$(x,y,z)$. Le Floch~\cite{lefloch}, following the methods of the Equational Theories
Project~\cite{etp}, determined all implications among the resulting \NLaws{} laws
with Prover9 and Mace4~\cite{mccune}. He states that every non-implication has a
finite countermodel, without giving their orders.

We measure, for each $k$, the part of this structure witnessed by quasigroups of
order at most $k$: non-implications, separation of classes, and closed sets
realised as intersections of signatures. Le Floch's result is the ground truth and
is not revised.

\section{Definitions}

A quasigroup of order $n$ is a Latin square on $\{0,\dots,n-1\}$. Following
\cite{lefloch}, $a \ldiv b$ is the $c$ with $c * a = b$, and $a \rdiv b$ is the $c$
with $b * c = a$. For a quasigroup $Q$, $\operatorname{sig}(Q)$ is the set of laws
$Q$ satisfies. $Q$ witnesses $A \not\Rightarrow B$ if $A \in \operatorname{sig}(Q)$
and $B \notin \operatorname{sig}(Q)$.

\begin{definition}\label{def:floor}
The floor at order $k$ is the set of quasigroups of order at most $k$ up to
isomorphism. A non-implication is witnessed at order $k$ if some quasigroup of the
floor witnesses it. A closed set $V$ is realised at order $k$ if
$V = \bigcap \{\operatorname{sig}(Q) : V \subseteq \operatorname{sig}(Q),\ |Q| \le k\}$,
with the empty intersection equal to the set of all laws.
\end{definition}

Every signature is closed under consequence, including consequences of conjunctions,
hence one of the \NVarieties{} closed sets listed by Le Floch. Equivalent laws have equal signatures, so the partition of laws by signature
over a floor is coarser than or equal to the partition into classes.

\begin{lemma}\label{lem:local}
Whether $A \not\Rightarrow B$ is witnessed at order $k$ depends only on the models of
$A$ of order at most $k$. Whether $V$ is realised at order $k$ depends only on the
models of $V$ of order at most $k$, and a closed set realised at order $k$ is realised
at every larger order.
\end{lemma}
\begin{proof}
The first two claims follow from the definitions of witness and realisation
(\cref{def:floor}). For the third, adding models of $V$ can only shrink the
intersection, which always contains $V$.
\end{proof}

\section{Method}

\paragraph{Enumeration.}
Quasigroups of order $n \le 6$ are enumerated by backtracking over Latin squares,
keeping the row-major lexicographically least table of each isomorphism class. Two
necessary conditions for leastness prune the search: $T_{00} \le 1$, and row 0 is
least among its conjugates by permutations fixing 0. Each law is checked on the
$n^3$ assignments, stopping at the first failure. Order 6 takes \EnumSecsSix{}\,s to
enumerate and \SigSecsSix{}\,s to evaluate on \Threads{} threads of an \CPU.

\paragraph{Verification.}
(i) The class counts \IsoOne, \IsoTwo, \IsoThree, \IsoFour, \IsoFive, \IsoSix{} match
OEIS \oeis{A057991}~\cite{oeis}, and $\sum n!/|\mathrm{Aut}|$ equals the number of
Latin squares (\oeis{A002860}) at each order. (ii) No floor quasigroup refutes a
listed implication (\FalseRefSix{} cases), and every signature is a listed closed
set. (iii) The counts of groups, abelian groups, commutative, semisymmetric, totally
symmetric and idempotent quasigroups agree with OEIS (\oeis{A000001},
\oeis{A000688}, \oeis{A057992}, \oeis{A076017}, \oeis{A076019}, diagonal of
\oeis{A058175}):
\GThreeAllMatch{} \GThreeCells{} values match. (iv) The six parastrophes map the law
list to itself; the predicted permutation of signatures was checked on every Latin
square of order at most \GFourMaxOrder{} (\GFourChecks{} bit comparisons,
\GFourMismatches{} mismatches). (v) A separate script, with its own law parser and
evaluator and reading Le Floch's files directly, re-checks \IndepCoversOk{} covers
below and finds a witness of the stated least order for all \IndepModelPairs{} residual
pairs of the Results section. SAT encodings solved with kissat confirm the least
orders: \Sch{38} and \Sch{79} have no model of order 7, no model of \Sch{23} of order 7
or 8 violates \Sch{5} (which implies \Sch{29}, \Sch{93} and \Sch{183}), and witnesses
exist at orders 8 and 9, re-checked after decoding. Two queries on associativity and
commutativity at orders 5 and 6 test the encoding. \SatAllOk{} \SatQueries{} queries
give the expected answer, each within \SatMaxSeconds{}\,s.

\paragraph{Beyond order 6.}
By \cref{lem:local} it suffices to enumerate, at order $k$, the models of each
closed set not realised at order $k-1$. This also covers witnessing: if
$A \not\Rightarrow B$ is not witnessed at order $k-1$, then $B$ lies in every signature
of order at most $k-1$ containing $A$, so the closed set generated by $A$ is not
realised at order $k-1$ and its models of order $k$ are enumerated. Separation of
classes follows from witnessing. The search fills cells row-major and prunes as soon
as a required law fails on defined entries, with the same symmetry breaking and
leastness test. With no laws, and with each of seven sample laws, it returns exactly
the enumerated isomorphism classes at every order at most 5. As an external check,
\Sch{5} is equivalent to semisymmetry~\cite{lefloch} and implies \Sch{23}; filtering
the \Sch{23} runs by \Sch{5} gives \SemiSeven, \SemiEight{} and \SemiNine{}
quasigroups at orders 7, 8 and 9, the values of \oeis{A076017}. The other runs have
no such external count.

\paragraph{Cover.}
Minimum sets of quasigroups witnessing a given set of non-implications are computed
as set-cover integer programs with HiGHS~\cite{highs}.

\section{Results}

\Cref{tab:results} gives the cumulative counts.

\begin{table}[ht]
\centering
\begin{tabular}{rrrrr}
\toprule
$k$ & class non-impl.\ witnessed & share & classes separated & closed sets realised\\
\midrule
3 & \ClassWitThree & \ClassSatThree & \ClassesSepThree & \VarThree\\
4 & \ClassWitFour & \ClassSatFour & \ClassesSepFour & \VarFour\\
5 & \ClassWitFive & \ClassSatFive & \ClassesSepFive & \VarFive\\
6 & \ClassWitSix & \ClassSatSix & \ClassesSepSix & \VarSix\\
\midrule
7 & \ClassWitSeven & \ClassSatSeven & \ClassesSepSeven & \VarSeven\\
8 & \ClassWitEight & \ClassSatEight & \ClassesSepEight & \VarEight\\
9 & \ClassWitNine & \ClassSatNine & \ClassesSepNine & \VarNine\\
\bottomrule
\end{tabular}
\caption{Cumulative results for the floor at order $k$. Denominators: \ClassNonImp{}
non-implications between classes, \NClasses{} classes, \NVarieties{} closed sets.
Orders 1 and 2 witness nothing. A class is separated if no other class has the same
signatures. Rows 7--9 come from the search justified by \cref{lem:local} and are
exact up to order \CompleteUpTo.}
\label{tab:results}
\end{table}

\paragraph{Small orders.}
Counted over laws, the floor at order 4 witnesses \LawSatFour{} of the \LawNonImp{}
non-implications, and the floor at order 6 \LawSatSix. The unique quasigroup of
order 2 satisfies all laws. $\mathbb{Z}_3$ satisfies \ZThreeLaws{} laws and alone
witnesses \ZThreeShare{} of the law non-implications.

\paragraph{Residual.}
At order 6, \UnrefSix{} class non-implications are unwitnessed, all with hypothesis
\Sch{23}, \Sch{38} or \Sch{79}. Their least witnesses have order \WitThirtyEight{}
(\Sch{38}), \WitSeventyNine{} (\Sch{79}) and \WitTwentyThree{} (\Sch{23}); Mace4 agrees
for \MaceAgree{} of the \ResidualPairs{} pairs. Small quasigroups miss these pairs
because the hypotheses have few models; for $n = 1, 2, \dots$ the model counts are
\Sch{7}: \CountsSeven; \Sch{38}: \CountsThirtyEight; \Sch{5}: \CountsFive;
\Sch{23}: \CountsTwentyThree. \Sch{7} implies \Sch{38}, so equal counts mean equal sets
of models: up to order \SevenThirtyEightSameUpTo{} no quasigroup separates them.

\paragraph{Closed sets.}
Least realising order of the \NGenerators{} closed sets not realised at order 6:
\VarLeastHist. Not realised at order at most \LastOrder: \MissingLastList.

\paragraph{Covers.}
The \CoverUniverse{} class non-implications witnessed at order at most 6 are
witnessed by \CoverILP{} quasigroups (\CoverProfile), optimal among quasigroups of
order at most 6; greedy uses \CoverGreedy{} and the LP bound is \CoverLP. For all
\ClassNonImp{} class non-implications, the optimum among known signatures is
\FullCover. An ILP over all \NVarieties{} closed sets as hypothetical signatures gives
the lower bound \FullCoverLB, valid for quasigroups of any order provided Le Floch's
list of closed sets is complete.

\paragraph{Symmetry.}
Under the parastrophes the \NLaws{} laws form \LawOrbits{} orbits, the classes
\ClassOrbits, the closed sets \VarOrbits{} and the ordered class pairs
\ClassPairOrbits.

\paragraph{Model counts.}
Counts of models by order for all classes and closed sets are in the repository. Of
the \OeisSeqs{} distinct class count sequences for orders 1--6, searching OEIS for the
terms together with ``quasigroup'', ``latin'' or ``loop'' finds \OeisMatched; for
\OeisNone{} of them these searches return nothing.

\section{Limitations}

The ground truth is Le Floch's computation and is not re-proved here. In particular,
the lower bound for covers assumes that his \NVarieties{} closed sets are all the
subsets of laws closed under consequence. Least witness orders beyond 6 are confirmed
by SAT; the realisation of closed sets at orders 7--9 rests on the completeness of our
search, checked as described. Mace4 is used only for comparison.
An OEIS match of six terms of a model-count sequence is not a proof. No result is
formalised in Lean.

\section*{Data and code}
\repo{} (Apache-2.0). Le Floch's ancillary files are redistributed under
\href{https://creativecommons.org/licenses/by/4.0/}{CC BY 4.0}.

\section*{Use of AI}
Code, analysis and text were produced with Claude (Anthropic) under the author's
direction. All numbers are generated by scripts in the repository.

\begin{thebibliography}{9}
\bibitem{etp} M.~Bolan, J.~Breitner, J.~Brox, et al., \emph{The Equational Theories
Project: Advancing Collaborative Mathematical Research at Scale}, \arxiv{2512.07087},
2025.
\bibitem{highs} Q.~Huangfu and J.~A.~J. Hall, \emph{Parallelizing the dual revised
simplex method}, Mathematical Programming Computation \textbf{10}(1) (2018),
119--142, \doi{10.1007/s12532-017-0130-5}.
\bibitem{lefloch} B.~Le Floch, \emph{Implication semilattice of 990 quasigroup
equational laws}, \arxiv{2603.29909}, 2026.
\bibitem{mccune} W.~McCune, \emph{Prover9 and Mace4},
\url{https://www.cs.unm.edu/~mccune/prover9/}, 2005--2010.
\bibitem{oeis} OEIS Foundation Inc., \emph{The On-Line Encyclopedia of Integer
Sequences}, \url{https://oeis.org}.
\bibitem{schroeder} E.~Schröder, \emph{Vorlesungen über die Algebra der Logik (exakte
Logik)}, vol.~1, Teubner, Leipzig, 1890.
\href{https://archive.org/details/bub_gb_P95LAAAAYAAJ}{Scan at archive.org}.
\end{thebibliography}

\end{document}
